\section{Conclusion}

This project demonstrates a complete and modern introductory HEP analysis workflow using LHCb CERN Open Data. Starting from ROOT ntuples, the analysis reads selected branches with \texttt{uproot}, processes candidates with \texttt{pandas} and \texttt{numpy}, reconstructs invariant masses, applies physics-motivated selections, generates publication-style plots, trains a sideband-labelled Random Forest classifier, performs charge-separated invariant-mass fits, computes a raw charge asymmetry, and visualises three-body phase space with Dalitz plots.

The invariant-mass reconstruction shows a clear B-candidate peak after the baseline selection. The selected mass spectrum is described with a double-Gaussian signal model and an exponential background model, providing a realistic detector-level signal extraction framework. The signal/background classification study demonstrates how sideband labels can be combined with standard machine-learning tools, while the ROC curve and feature-importance plot provide useful performance diagnostics. The CP-asymmetry study separates \Bplus{} and \Bminus{} candidates, performs independent invariant-mass fits, and extracts a raw integrated asymmetry together with statistical and systematic uncertainty estimates. Toy-MC pseudo-experiments are used to validate the statistical behaviour of the fitting procedure and uncertainty propagation. The Dalitz analysis extends the project toward multibody-decay phase-space studies by reconstructing pairwise invariant-mass-squared variables.

The analysis remains intentionally simplified and does not constitute a precision flavour-physics measurement. Several improvements would still be required for a full experimental analysis:
\begin{itemize}
  \item train stronger classifiers such as XGBoost or gradient-boosted decision trees, with cross-validation and robust overtraining checks;
  \item include detector effects, PID calibration, trigger efficiencies, and reconstruction asymmetries;
  \item use simulation and control samples to validate selection efficiencies and background shapes;
  \item combine MagnetUp and MagnetDown samples to reduce charge-dependent detector biases;
  \item implement additional accepted fit-model variations, such as Crystal Ball signal models or alternative background parameterisations;
  \item perform local CP-asymmetry studies across Dalitz-plot regions.
\end{itemize}

Despite these limitations, the project captures several essential components of a realistic experimental-particle-physics analysis: transparent event reconstruction, documented selection choices, modular code, reproducible figures, signal extraction through invariant-mass fitting, uncertainty propagation, toy-MC validation, and explicit treatment of systematic effects. It therefore provides a strong foundation for more advanced studies in heavy-flavour physics and for future work in an experimental PhD environment.